\documentclass{strategynet-report}

\hypersetup{
  pdftitle={Volt: Strategy Search},
  pdfauthor={research@strategynet.ai}
}

\lstdefinestyle{volt}{
  style=sncode
}

\newcommand{\VOLT}{\textsf{Volt}}
\newcommand{\COG}{\texttt{cog}}
\newcommand{\E}{\mathbb{E}}

\tikzset{
  discoverybox/.style={
    draw=SNLine,
    fill=SNPaper,
    rounded corners=2pt,
    text width=26mm,
    minimum height=18mm,
    align=center,
    inner sep=2.5mm,
    font=\sffamily\fontsize{8}{9.5}\selectfont
  },
  discoverylearn/.style={
    discoverybox,
    draw=SNSignalDark,
    fill=SNSignal!6
  },
  discoveryflow/.style={
    -{Latex[length=2mm]},
    draw=SNInk,
    line width=0.75pt
  },
  discoveryfeedback/.style={
    -{Latex[length=2mm]},
    draw=SNSignalDark,
    dashed,
    line width=0.75pt
  }
}

\title{\VOLT{}: Strategy Search}
\SNSubtitle{Finding useful portfolios without evaluating every candidate}
\author{\href{mailto:research@strategynet.ai}{research@strategynet.ai}}
\SNDocumentType{Research architecture}
\SNVersion{Version 0.2}
\SNDate{July 2026}
\SNStatus{Open-source working paper}
\SNShortTitle{Volt: Strategy Search}

\begin{document}
\maketitle

\begin{abstract}
An option chain can produce more valid portfolios than a researcher can price,
simulate, and review.  Many candidates are duplicates or slight variations;
others look attractive only because their features used stale information or
their labels ignored executable prices.  Exhaustive evaluation wastes a
limited research budget, while an unconstrained ranker can learn the wrong
lesson very efficiently.

This paper describes a search process that spends evaluation effort on
distinct, legal, point-in-time candidates.  Each candidate becomes a
versioned \VOLT{} graph linked to its exact portfolio identity, bounded
actions, and deterministic simulation evidence.  The proposed statistical
layer learns over those graphs and adds dataset construction, candidate
scoring, active acquisition, baseline comparisons, quote-backed labels, and
chronological evaluation.  The public repository ships synthetic fixtures
and semantic reduction results; it does not ship a trained ranker or make an
alpha claim.  The initial \COG{} compiler, compiler-bound graph dataset, and
structural legal mask are implemented; account-aware margin and capital policy
remain integration work.
\end{abstract}

\SNtableofcontents

\part{Research Design}

\section{Discovery task}

At a decision time \(t\), the research system receives a market and account
state together with the contracts that were listed and observable at that
time.  The implemented \texttt{CandidateGrammar} resolves typed actions
against the caller-supplied listed universe and enforces structural and
exercise-authority rules.  Caller policy must still apply liquidity, capital,
margin, and account constraints.  The implemented \COG{} compiler coordinates
the structural decisions and their diagnostics; the discovery model should see
candidates only after every deterministic decision has been recorded.

The initial task is bounded.  Given an admissible set
\(A_t\), estimate the outcome and risk of each candidate and rank the set under
a declared objective.  A candidate may be a complete portfolio constructed
from an empty account or a one-step transformation of existing inventory.
Either formulation yields an ordinary supervised sample once the simulator or
historical replay has supplied the outcome.  Exhaustive evaluation remains
possible on small universes, which gives the experiment a known oracle and
makes search efficiency measurable.

This separation keeps the research question precise.  Contract validity,
exercise authority, and settlement remain Volt runtime concerns; account
admissibility remains an application concern.  The model is
judged on how effectively it orders legal alternatives, how much expensive
evaluation it saves, and whether its ranking remains useful after spreads,
fees, and out-of-time market changes.

\section{Artifact contract}

The proposed research interface exposes a sample
\[
  d_t=(G_t,A_t,M_t,y_t,p_t),
\]
where \(G_t\) is the compiled graph, \(A_t\) contains feasible candidates,
\(M_t\) is the legal mask, \(y_t\) carries outcome and risk labels, and \(p_t\)
records provenance.  Every element is versioned.  The model may replace its
graph schema, normalizers, features, and objective without changing the
authoritative portfolio or the events that produced its label.  The current
\texttt{PortfolioGraphSample} implements \(G_t\), exact and payoff hashes,
and an optional compact target.  Volt does not yet serialize the complete
five-part sample as one artifact.

\begin{center}
\begin{tabularx}{\textwidth}{@{}p{0.13\textwidth}X X@{}}
\toprule
Field & Financial content & Experimental use \\
\midrule
\(G_t\) & portfolio, exact option positions, underlying, expiry, and typed relations &
model input and attribution anchor \\
\(A_t\) & resolved candidate portfolios or transformations &
objects to score under one acquisition budget \\
\(M_t\) & authority, lifecycle, liquidity, and resource decisions &
exclusion before ranking or normalization \\
\(y_t\) & P\&L, drawdown, tail risk, turnover, costs, and scenario outcomes &
regression, ranking, calibration, or Pareto targets \\
\(p_t\) & timestamps, sources, hashes, assumptions, and split membership &
replay, leakage checks, and result identity \\
\bottomrule
\end{tabularx}
\end{center}

The experiment manifest fixes the observation clock, prediction horizon,
eligible grammar, execution models, scenario distribution, random seeds,
objective, acquisition budget, and chronological split before fitting begins.
Normalizers and learned categorical vocabularies use training data alone.
Validation data select configurations; the test split remains sealed until
the configuration and simulator budget have been frozen.

\section{Candidate pool}

The public fixtures use six synthetic same-expiry contracts.  In compiled
mode, the \texttt{graph} example constructs and compiles one versioned source
program for each quantity vector in \(\{-1,0,1\}\), producing all
\(3^6=729\) compiler-bound graph samples.  Each sample carries source and
artifact hashes, a deterministic identifier, an exact portfolio hash, an
optional terminal-payoff hash, typed nodes and edges, and a synthetic target.

The \texttt{zap} example supplies two semantic controls.  Two ordered actions
produce 144 raw sequences but 73 exact portfolios.  Independently, quantities
in \(\{-2,-1,0,1,2\}\) across six contracts produce 15,625 exact portfolios
and 13,321 terminal-payoff classes.  These fixtures test identity,
canonicalization, and export plumbing; they contain no historical market
cohort, trained model, liquidity filter, or economic claim.

As the universe expands, exhaustive labels become too costly.  Candidate
generation and evaluation must then be treated as separate budgets.  The
\COG{} compiler or caller application may enumerate or sample a large valid
pool cheaply, while the acquisition policy chooses the smaller subset that
receives scenario simulation or historical replay.

\part{Models}

\section{Graph view}

\subsection{Schema}

The implemented graph is derived from a \texttt{CanonicalPortfolio} plus a
validated \texttt{GraphContext}.  Its node types are \texttt{Portfolio},
\texttt{Underlying}, \texttt{Expiry}, and \texttt{OptionPosition}; its edge
types are \texttt{PositionIn}, \texttt{WrittenOn}, and \texttt{ExpiresAt}.
Every sample carries both exact-portfolio and terminal-payoff hashes.

Position features include contract identity, right, settlement, exercise
style, strike, declared mark, signed quantity, multiplier, first exercise
time, and discrete opportunity count.  Context supplies one underlying,
currency, valuation time, common expiry, reference spot, and a mark for every
contract.  Account nodes, Greeks, spread, liquidity, lifecycle state, market
events, and multi-expiry graphs are not part of schema version
\texttt{volt.portfolio\_graph.v1}.

The current view accommodates variable leg counts while remaining insensitive
to canonical input order.  Future experiments may derive richer graph schemas
without rewriting contract semantics, but each schema requires its own
version and controls.  A set encoder and conventional engineered features
remain essential comparisons \cite{relational}.

\subsection{Encoding}

Type-specific input transforms map nodes into a common latent dimension.
Relation-specific message passing produces contextual node states, and a
permutation-invariant readout gives a portfolio representation \(h_t\).
Candidate \(a\) receives its own representation \(e(a)\), either from a
separately compiled graph or from the proposed graph edit.  One useful score
is
\[
  Q_\theta(G_t,a)
    = f_\theta\!\left(h_t,e(a),h_t\odot e(a)\right).
\]
Companion heads may estimate an outcome distribution, named risk measures,
and a behavioral embedding.  The legal mask is applied before ranking or
probability normalization, so an invalid action cannot acquire probability
through a numerical accident.

\section{Training objective}

No single scalar captures every use of an option portfolio.  A research
dataset may therefore derive a vector
\[
\begin{aligned}
  y=(&\text{net P\&L},\text{drawdown},\text{CVaR},\text{turnover},\\
     &\text{capital use},\text{exercise loss},\text{scenario vector}),
\end{aligned}
\]
from \texttt{SimulationResult} and additional named risk models.  The current
result directly records fills, journals, pathwise and terminal P\&L, leg
attribution, final state, and reproducibility hashes; it does not calculate
drawdown, CVaR, or capital use automatically.  The manifest states how the
experiment derives and uses any added targets.  A fitting loss may
combine outcome regression, pairwise or listwise ranking, distribution
calibration, and auxiliary risk targets:
\[
  \mathcal L
    =\lambda_v\mathcal L_{\mathrm{value}}
    +\lambda_r\mathcal L_{\mathrm{rank}}
    +\lambda_c\mathcal L_{\mathrm{calibration}}
    +\lambda_a\mathcal L_{\mathrm{aux}}.
\]
Coefficients, target transforms, and any scalar utility belong to the
experiment record.  Changing them creates a new study while leaving candidate
identity and financial labels intact.

\section{Baselines}

Graph models must justify their additional machinery against simpler
alternatives.  Every comparison includes deterministic trading heuristics,
linear or tree models on engineered features, a permutation-invariant set
encoder, and random or unguided acquisition.  All models receive identical
point-in-time information, candidate pools, labels, and evaluation budgets.
Wall-clock measurements accompany simulator-call counts because a model that
saves evaluations but consumes much more compute may still be unattractive.

Deep Sets provides the closest structural control: it handles variable-sized,
order-invariant collections without relation-specific message passing
\cite{deepsets}.  GraphGym-style sweeps can standardize optimization and
architecture comparisons \cite{graphgym}; PyTorch Geometric supplies the
implementation substrate \cite{pyg}.  The Volt manifest continues to own
financial objectives, data splits, and evidence identities.

\section{Acquisition loop}

\begin{center}
\begin{tikzpicture}[node distance=4mm]
  \node[discoverybox] (data) {\textbf{Volt samples}\\graphs, masks, labels,
    evidence};
  \node[discoverylearn,right=of data] (fit) {\textbf{Fit}\\value, rank, risk,
    embedding};
  \node[discoverybox,right=of fit] (rank) {\textbf{Rank pool}\\score legal,
    unlabeled candidates};
  \node[discoverylearn,right=of rank] (acquire) {\textbf{Acquire}\\top,
    uncertain, and random};
  \node[discoverybox,right=of acquire] (simulate) {\textbf{Evaluate}\\simulate,
    validate, append};
  \draw[discoveryflow] (data) -- (fit);
  \draw[discoveryflow] (fit) -- (rank);
  \draw[discoveryflow] (rank) -- (acquire);
  \draw[discoveryflow] (acquire) -- (simulate);
  \draw[discoveryfeedback] (simulate.south) -- ++(0,-5mm) -|
    (data.south);
\end{tikzpicture}
\begingroup
\captionsetup{hypcap=false}
\captionof{figure}{Supervised discovery under a fixed acquisition budget.
The feedback path carries immutable simulator evidence into a new dataset
version; it never alters an existing label.}
\label{fig:acquisition-loop}
\endgroup
\end{center}
\par\medskip

An initial labeled cohort fits the value and ranking heads.  The model then
scores an unlabeled pool, after which the acquisition rule selects candidates
for expensive evaluation.  High predicted utility alone is insufficient:
early errors can concentrate the dataset around an accidental pattern.  The
batch therefore combines promising candidates with uncertainty or diversity
coverage and a compulsory random allocation.

Each acquisition records the candidate universe, model checkpoint, score,
selection reason, and probability where available.  Once evaluated, the new
evidence is appended under a new dataset manifest and the model is retrained.
Learning curves plot result quality against both simulator calls and elapsed
compute.  The experiment stops at its preregistered budget rather than when a
visually attractive backtest appears.

\part{Evaluation}

\section{Quote-backed labels}

Marks and fills remain separate.  \texttt{ScenarioOptionPrice} records a
declared mark and optional bid and ask.  \texttt{DeclaredMark} uses the mark as
a synthetic or proxy fill; \texttt{CrossAtTouch} buys at the ask and sells at
the bid, failing closed when the required side is absent.  Neither convention
establishes that an order executed.

The public repository does not ship historical quote-backed rankings.  A
future ranking experiment should compare declared marks, quote midpoints, and
executable-side conventions, record displayed-size and fill limitations, and
report sensitivity to the spread rather than importing unpublished
correlations into this paper.

\section{Splits and metrics}

No train/validation/test market cohort is included in Volt.  A future study
must use chronological splits and, where repeated issuers are present, an
issuer-aware split so company behavior cannot cross freely between training
and evaluation.

Primary measurements include:
\begin{itemize}[nosep]
  \item precision and normalized discounted cumulative gain at \(k\);
  \item recovery of exhaustive-search top-decile candidates;
  \item rank correlation and top-\(k\) overlap across execution assumptions;
  \item distribution calibration and tail-risk error;
  \item diversity and canonical uniqueness of acquired candidates;
  \item out-of-time utility after declared costs;
  \item simulator calls and elapsed compute required to reach a target score.
\end{itemize}
Invalid-action rate should remain zero after masking.  Volt now validates the
alignment and integrity of the compiler's structural mask; a production study
must compose it with caller-owned liquidity, margin, and account policy.

\section{Research sequence}

\subsection{Canonical search}

The bounded enumeration supplies ground truth for raw constructions, exact
portfolios, payoff classes, dominance, and clustering.  It supports direct
measurement of reduction ratios and lets the ranker be compared with
exhaustive search before acquisition becomes necessary.

\subsection{Event portfolios}

An earnings study expands the cohort to liquid U.S. names observed immediately
before scheduled releases.  Delta buckets and expirations remain narrow enough
to preserve an inspectable universe.  Post-event labels use executable quote
conventions, and time and issuer splits limit memorization.

\subsection{Volatility dynamics}

Further studies separate fixed-contract changes from relative-coordinate
changes in strike and maturity.  Term structure, forward variance, skew carry,
and spot--volatility correlation are estimated by regime and observation
horizon.  Their features must be constructed entirely from information
available at the ranking clock.

\subsection{Conformance}

The current Rust tests verify exact canonicalization, terminal hinge payoff,
exercise and assignment consequences, deterministic replay, quote authority,
scenario fills, graph export, and candidate branching.  Future macro support
for \emph{verticals}, \emph{straddles}, \emph{strangles},
\emph{butterflies}, \emph{iron condors}, and \emph{calendars}
will require independent expansion and ordering fixtures.  Such controls
distinguish a model failure from a financial representation error.

\section{Evidence and limits}

Every prediction links to the graph and checkpoint that produced it.  Graph
attributions and feature ablations describe the model; leg and scenario
attribution reconcile the financial outcome; the Volt journal establishes the
contracts, observations, and transitions.  Keeping those records separate
prevents a persuasive model explanation from substituting for accounting
evidence.

The grammar, market cohort, simulation model, and objective place hard limits
on discovery.  A learner may exploit an optimistic fill convention, a data
revision, a simulator defect, or a narrow market regime as readily as it
exploits a durable relationship.  Candidates selected for further study
therefore face adversarial scenarios, alternative costs, capacity analysis,
and untouched time periods before their apparent utility receives an economic
interpretation.

The current synthetic fixtures verify infrastructure only.  Their objective
and scenario weights are declared synthetic inputs, and their
cross-at-touch behavior is a modeled fill rule.  The next useful research
result is an out-of-time acquisition curve on a public, point-in-time cohort,
compared with the full baseline suite at equal evaluation budgets.

\section{Conclusion}

Volt supplies a stable financial search space; the discovery system learns how
to traverse it economically.  This study keeps the decision static so
candidate identity, labels, rankings, and acquisition cost can be audited
against exhaustive ground truth.  Once that loop is reliable, the same graph
encoder may support sequential control, but trajectory generation, delayed
reward, and rollout allocation belong to a different experimental problem.

\begin{thebibliography}{9}

\bibitem{deepsets}
M. Zaheer, S. Kottur, S. Ravanbakhsh, B. Póczos, R. Salakhutdinov, and
A. Smola,
\emph{Deep Sets},
Advances in Neural Information Processing Systems, 2017,
\url{https://arxiv.org/abs/1703.06114}.

\bibitem{relational}
P. Battaglia et al.,
\emph{Relational Inductive Biases, Deep Learning, and Graph Networks},
2018,
\url{https://arxiv.org/abs/1806.01261}.

\bibitem{pyg}
M. Fey and J. E. Lenssen,
\emph{Fast Graph Representation Learning with PyTorch Geometric},
ICLR Workshop on Representation Learning on Graphs and Manifolds, 2019,
\url{https://arxiv.org/abs/1903.02428}.

\bibitem{graphgym}
J. You, R. Ying, and J. Leskovec,
\emph{Design Space for Graph Neural Networks},
Advances in Neural Information Processing Systems, 2020,
\url{https://arxiv.org/abs/2011.08843}.

\end{thebibliography}

\end{document}
